\documentclass[11pt,a4paper]{article}
\usepackage[T1]{fontenc}
\usepackage{lmodern,amsmath,amssymb,amsthm,booktabs,microtype}
\usepackage[margin=22mm]{geometry}
\usepackage[hidelinks]{hyperref}
\newtheorem{theorem}{Theorem}
\newtheorem{lemma}{Lemma}
\newcommand{\SC}{\operatorname{SC}}
\newcommand{\comp}{\operatorname{comp}}
\newcommand{\Sn}{\mathfrak S_n}
\title{A length-nine gap in the fertility spectra\protect\\of consecutive-pattern stack-sorting maps}
\author{Aurantium}
\date{14 October 2026}
\begin{document}
\maketitle
\begin{abstract}
For each of the single forbidden patterns $213$ and $231$, the fertility
spectrum of the consecutive-pattern-avoiding stack-sorting map on
$\mathfrak S_9$ is
$\{0,128\}\cup(\{1,\ldots,64\}\setminus\{55\})$.
An archived complete enumeration visits all $9!=362\,880$ inputs for
$\SC_{213}$ and records the full histogram, including $265\,602$ outputs
with zero fertility. A step-by-step complement conjugacy gives the same
histogram for $\SC_{231}$. The missing value $55\le2^{9-3}$ contradicts
the fixed-length inclusion in Kemeklis's Conjecture~4.1. It does not
contradict the theorem that every positive integer is a fertility when
the length may vary. We specify the algorithm, its correspondence with
the original definition, and the exact scope of the recorded execution.
\end{abstract}

\section{The fixed-length statement}
Write $\Sn$ for the permutations of $[n]=\{1,\ldots,n\}$. Fix a
\emph{single} pattern $\sigma\in\mathfrak S_3$. The map $\SC_\sigma$
reads its input from left to right and maintains a stack read from
\emph{top to bottom}. With next input $x$, test the consecutive top
triple that would result from pushing $x$. If its relative order is
$\sigma$, pop the current top to the output and retry the same $x$.
Otherwise push $x$ and advance the input. After the input is consumed,
pop the remaining stack to the output in top-to-bottom order. This is
the rule in \cite[p.~1]{K}.

For an output $\pi\in\Sn$, define
\[
 f_{\sigma,n}(\pi)=\#\{\tau\in\Sn:\SC_\sigma(\tau)=\pi\},
 \qquad A_n^\sigma=\{f_{\sigma,n}(\pi):\pi\in\Sn\}.
\]
Zero is included when some output is not attained.

Kemeklis's Conjecture~4.1 \cite[p.~11]{K} proposes the inclusion
\begin{equation}\label{eq:conjecture}
 \{1,2,\ldots,2^{n-3}\}\subseteq A_n^\sigma,
 \qquad \sigma\in\{213,231\}.
\end{equation}
For integer $n\ge3$, this asks that every integer in the displayed
range occur at that same length. The printed source uses the notation
$\sigma=\{213,231\}$ and the symbol $\subset$. In the context of the
single-pattern maps defined there, we read this as the two separate
choices $\sigma=213$ and $\sigma=231$, not a machine simultaneously
avoiding both patterns. Our counterexample uses $n=9$; it violates even
the weak inclusion in \eqref{eq:conjecture}, independently of whether
the source's subset symbol is intended to require strict inclusion.

\begin{theorem}[Recorded exhaustive computation and symmetry]\label{thm:main}
For the two single-pattern maps,
\begin{equation}\label{eq:spectrum}
 A_9^{213}=A_9^{231}
   =\{0,128\}\cup\bigl(\{1,\ldots,64\}\setminus\{55\}\bigr).
\end{equation}
In particular, no length-nine output of either map has fertility $55$,
so \eqref{eq:conjecture} fails at $n=9$.
\end{theorem}

The theorem has a finite computational component: it depends on the
specified program and its completed execution. The symmetry component
is proved below. This note does not claim a hand classification of all
fibers or proof-assistant certification.

\section{Why the program implements the defined map}
For three distinct numbers $a,b,c$, standardization replaces each
entry by one plus the number of smaller entries. For a prospective
triple $(x,s_0,s_1)$, with $s_0$ the stack top and $s_1$ immediately
below it, the order is $213$ exactly when
\begin{equation}\label{eq:blocked}
 s_0<x<s_1.
\end{equation}
The recorded independent implementation computes all three ranks
directly, rather than substituting a specialized inequality test.
Its complete stack function is:
\begin{quote}\small
\begin{verbatim}
def sc(p):
    stack=deque(); output=[]
    for x in p:
        while len(stack)>=2:
            triple=(x,stack[0],stack[1])
            order=tuple(1+sum(b<a for b in triple)
                        for a in triple)
            if order!=(2,1,3): break
            output.append(stack.popleft())
        stack.appendleft(x)
    return tuple(output)+tuple(stack)
\end{verbatim}
\end{quote}
Here the left end of the deque is the top. The line wrapping of the
rank expression is only typographical; the distributed program preserves
the original mathematical function and enumeration logic.

\begin{lemma}\label{lem:implementation}
For every $p\in\Sn$, this function terminates and returns $\SC_{213}(p)$,
with each input label appearing once in the output.
\end{lemma}
\begin{proof}
Initially the stack is empty and avoids the forbidden consecutive
pattern. Popping its top leaves a suffix of the top-to-bottom sequence,
so cannot create a new consecutive triple. A push creates only one
new triple, namely $(x,s_0,s_1)$, when two entries were already in the
stack. Thus testing only that triple suffices to maintain the invariant.
If the stack has fewer than two entries, the push cannot create a
length-three pattern.

The inner loop pops precisely when the prospective triple standardizes
to $(2,1,3)$. A pop does not change the outer-loop input $x$, so the
same input is retried against the updated stack. Breaking the loop,
or reducing the stack below length two, is followed by pushing $x$
at the top. This matches each decision in the defining rule.

For any fixed input entry, each unsuccessful trial removes a stack
entry, so the inner loop terminates. Every label is pushed once and
popped once, either during a trial or in the final drain. No label is
created, duplicated or discarded. The final deque is already listed
top-to-bottom, hence concatenating its tuple to the output is exactly
the final drain. The output is therefore a permutation of $[n]$, and
every step agrees with the original deterministic map.
\end{proof}

\section{Enumeration, accounting and the complete histogram}
The original computation iterates over
\texttt{permutations(range(1,10))}. The permutation iterator supplies
each ordering of the nine distinct labels exactly once, hence every
element of $\mathfrak S_9$. The loop increments an input counter and,
for each input $p$, increments \texttt{counts[sc(p)]} by one. Keys are
complete ordered output tuples, not abbreviated permutation ranks or
hash values. The counter uses exact Python integers; dictionary hash
collisions are resolved by tuple equality and do not merge distinct
outputs.

By Lemma~\ref{lem:implementation}, each loop iteration contributes once
to the fiber of its unique output. Thus at the end, for every attained
$\pi$, \texttt{counts[pi]} is exactly $f_{213,9}(\pi)$. There is no
restriction to inputs whose outputs are sorted. The original program
asserts that its input counter equals $9!$.

Define the full histogram
\[
 h_9(f)=\#\{\pi\in\mathfrak S_9:f_{213,9}(\pi)=f\},
 \qquad f\in\mathbb Z_{\ge0}.
\]
Applying a second counter to \texttt{counts.values()} gives its positive
bins. Because the output universe is all of $\mathfrak S_9$, the
unattained outputs must also be counted:
\begin{equation}\label{eq:zero}
 h_9(0)=9!-\#\{\pi:f_{213,9}(\pi)>0\}
       =362\,880-97\,278=265\,602.
\end{equation}
This step is essential: zero-fertility outputs do not appear among
the dictionary keys.

The archived result is transcribed in Table~\ref{tab:histogram}. The
completed original program also checked the exact accounting identities
\begin{equation}\label{eq:accounting}
 \sum_{f\ge0}h_9(f)=362\,880,
 \qquad \sum_{f\ge0}f\,h_9(f)=362\,880.
\end{equation}
The first counts all outputs, including zero fibers; the second
partitions all inputs by their output. These consistency checks do
not replace the coverage argument, but detect missing mass in the
recorded histogram. The asserted support was tested \emph{after} the
counts were computed; the expected set did not generate the counts.

\begin{table}[p]
\centering
\caption{Complete length-nine histogram for $\SC_{213}$. Each pair of
columns is a fertility $f$ and the number $h_9(f)$ of outputs. The
$f=55$ row explicitly displays zero; every unlisted $f$ also has count
zero. The same histogram holds for $\SC_{231}$ by complement.}
\label{tab:histogram}
\medskip
\begin{tabular}{rr@{\qquad}rr@{\qquad}rr}
\toprule
$f$&$h_9(f)$&$f$&$h_9(f)$&$f$&$h_9(f)$\\
\midrule
0 & 265\,602 & 22 & 164 & 44 & 7 \\
1 & 40\,163 & 23 & 58 & 45 & 2 \\
2 & 11\,353 & 24 & 98 & 46 & 5 \\
3 & 11\,534 & 25 & 75 & 47 & 5 \\
4 & 9\,915 & 26 & 182 & 48 & 10 \\
5 & 5\,736 & 27 & 69 & 49 & 4 \\
6 & 3\,595 & 28 & 79 & 50 & 1 \\
7 & 3\,684 & 29 & 37 & 51 & 1 \\
8 & 2\,039 & 30 & 55 & 52 & 4 \\
9 & 1\,297 & 31 & 91 & 53 & 3 \\
10 & 1\,308 & 32 & 60 & 54 & 4 \\
11 & 1\,470 & 33 & 7 & 55 & 0 \\
12 & 854 & 34 & 3 & 56 & 5 \\
13 & 569 & 35 & 3 & 57 & 9 \\
14 & 663 & 36 & 12 & 58 & 3 \\
15 & 667 & 37 & 5 & 59 & 2 \\
16 & 607 & 38 & 7 & 60 & 4 \\
17 & 164 & 39 & 5 & 61 & 1 \\
18 & 144 & 40 & 3 & 62 & 3 \\
19 & 182 & 41 & 2 & 63 & 6 \\
20 & 148 & 42 & 9 & 64 & 9 \\
21 & 69 & 43 & 4 & 128 & 1 \\
\bottomrule
\end{tabular}
\end{table}

The full census has $97\,278$ positive-fertility outputs. Its support
is exactly the right-hand side of \eqref{eq:spectrum}: there are
64 distinct positive fertility values, the largest is $128$, and $55$ is
the unique missing value in $\{1,\ldots,64\}$. These are statements
about $n=9$ only.

\section{Complement conjugacy, including empty fibers}
For a fixed length $n$, let $c_n$ act entrywise by $x\mapsto n+1-x$.
It is an involutive bijection of $\Sn$. Complement symmetry is already
stated in \cite[Theorem 3.5 and Corollary 3.6]{K}; we give the complete
operation-level argument to fix the orientation and the zero case.

\begin{lemma}\label{lem:complement}
For every $n\ge1$ and $\tau\in\Sn$,
\begin{equation}\label{eq:conjugacy}
 \SC_{231}(c_n(\tau))=c_n(\SC_{213}(\tau)).
\end{equation}
Consequently $f_{231,n}(c_n(\pi))=f_{213,n}(\pi)$ for every
$\pi\in\Sn$, including when either fiber is empty.
\end{lemma}
\begin{proof}
Run the two machines in parallel. Maintain the invariant that they
have consumed the same number of inputs, their pending inputs are
entrywise complements, and their stacks and outputs are entrywise
complements in the same top-to-bottom and left-to-right orders.
This holds initially with empty stacks and outputs.

For a triple of distinct entries, complementing values reverses all
value comparisons. Its rank tuple $(a,b,c)$ becomes $(4-a,4-b,4-c)$;
in particular $(2,1,3)$ becomes $(2,3,1)$. Therefore the prospective
push is forbidden in the first machine exactly when it is forbidden
in the second. If forbidden, both pop their current tops to their
outputs and neither advances the pending input; the invariant is
preserved. Otherwise both push their pending inputs and advance by
one, again preserving the invariant. A stack of size below two
admits the push on both sides. Once inputs are exhausted, their
top-to-bottom drains also preserve the invariant. Termination then
gives \eqref{eq:conjugacy}. This is value complement, with no reversal
of positions or stack orientation.

For a fixed $\pi$, the map $\tau\mapsto c_n(\tau)$ sends the fiber
of $\pi$ under $\SC_{213}$ into the fiber of $c_n(\pi)$ under
$\SC_{231}$. Applying the same complement again gives the inverse,
so the fibers are in bijection. The argument also proves that one
fiber is empty if and only if the other is empty. Finally the
output map $\pi\mapsto c_n(\pi)$ is bijective, so the entire fertility
histograms, not just their supports, are equal.
\end{proof}

\begin{proof}[Proof of Theorem~\ref{thm:main}]
The completed exhaustive computation described above gives
Table~\ref{tab:histogram} and hence $A_9^{213}$ as in
\eqref{eq:spectrum}. Lemma~\ref{lem:complement} transfers this
histogram to $\SC_{231}$ without a separate enumeration of that map.
Since $55\in\{1,\ldots,2^{9-3}\}$ and $h_9(55)=0$, the required
fixed-length inclusion fails for each map.
\end{proof}

\section{Recorded execution and trust boundary}
The original run was recorded on 13 October 2026. It started at
04:27:46.4455254 UTC and ended at 04:27:51.1014079 UTC.
The process exit code is $0$,
the stop reason is null, and stderr is empty. The program's internal
census timer records 4.4192823 seconds. The process monitor records
an observed peak working set of $31\,113\,216$ bytes, sampled every 50 ms;
this is the archived measurement, not a new measurement or a claim
about total system memory. The original receipt identifies an
interpreter executable but does not record its version, so none is
inferred here.

The original program first computed and saved the complete histogram,
then compared it with an author-supplied histogram. The recorded
stdout confirms that comparison passed. The result is supported by
the actual independent enumeration and its exit record, not merely
by agreement between two previously supplied data files. This note
does not claim a new execution of either program.

The data file \texttt{data/n9.recomputed.json} is preserved byte-for-byte
from that run. Its SHA-256 is
\begin{center}\small\ttfamily
ccf28642e4004512df0bd4054c0f98ce85d007442cc42effa17e697d2bbed882
\end{center}
The distributed stdout and empty stderr are likewise unchanged;
\texttt{evidence/EXECUTION.json} is a path-free summary binding the
original script, receipt, monitor, result and output hashes.

The portable program \texttt{code/reproduce\_n9.py} keeps the original
stack function, input iteration, counters, assertions and result
construction unchanged at the abstract-syntax-tree level. Its only
adaptations are a neutral header and stdout-only output, replacing
the old file-writing and path-specific comparison footer. This
\emph{packaged entry has not been executed}. Static preservation of
the mathematics does not turn its new file bytes into an already
executed artifact; the completed run remains bound to the original
script. Further details and hashes are in the reproducibility note.

This is an exhaustive computational counterexample, conditional on
the short implementation, Python's permutation and integer/container
semantics, and the archived execution record. It is not a formal
kernel certificate. The finite table and retained code permit direct
scrutiny without claiming new independent human peer review.

\section{What is and is not contradicted}
Kemeklis's Theorem~1.1 \cite[p.~2]{K} has different quantifiers:
for every positive integer $f$ and every $\sigma\in\mathfrak S_3$,
there exist a length $n$ and a permutation $\pi\in\Sn$ whose
fertility is $f$. That length may depend on $f$ and $\sigma$.
Absence of fertility $55$ at one fixed length therefore does not
contradict it. No length-ten calculation or witness is needed for
this distinction or for the refutation at length nine.

No minimal-counterexample-length claim is made. No simultaneous
avoidance machine is considered. This note neither rechecks larger
scans nor certifies worldwide publication priority.
Prepared with AI assistance.

\begin{thebibliography}{1}
\raggedright
\bibitem{K}
Jurgis Kemeklis, \emph{Fertility Numbers of Consecutive $S_3$
Pattern-Avoiding Stack-Sorting Maps},
\href{https://arxiv.org/pdf/2408.05378v2}{arXiv:2408.05378v2},
15 September 2024. Algorithm p.~1; Theorem~1.1 p.~2;
Theorem~3.5 and Corollary~3.6; Conjecture~4.1 p.~11.
\end{thebibliography}
\end{document}
